This technique has been brought into the 21st century with high-tech "Targeted Dream Incubation (TDI)" devices, such as the "Dormio" system developed at the MIT Media Lab \citep{haarhorowitz2020dormio}.

This device provides experimental proof for the theories discussed in Section 2. The Dormio glove-like system tracks the wearer's biosignals (muscle tone, heart rate, skin conductance) to detect the \textit{precise onset of the hypnagogic (N1) state}. At that moment, the system plays a pre-recorded audio cue (e.g., "Remember to think of a tree"). It then wakes the user, collects a "dream report," and allows them to drift back into the N1 state, repeating the process \citep{haarhorowitz2020dormio}.

The Dormio study validated the link between hypnagogia and creativity. Participants who napped with the "tree" incubation (TDI) performed \textit{significantly} more creatively on post-sleep creative tasks related to the "tree" topic. This was the first study to demonstrate that "dreaming of a specific topic during sleep onset is directly related to increased post-sleep creativity on that topic" \citep{haarhorowitz2020dormio}. The Dormio device, in effect, mechanizes the "dozing" state that Kekulé experienced by chance, targeting the specific "sweet spot" of N1 sleep for creative synthesis.